% !PRINTED 666
\weberpage{666}

%% [running head, not body text -- commented out 2026-08-07] Page furniture (running
%% head, folio, signature) is excluded from this edition by standing policy.
%% 666 \hfill \gesperrt{Twenty-second Chapter.} \hfill \S.\,183.

\begin{center}
\S.\,183.

\gesperrt{Proof of the second lemma for an odd $m$.}
\end{center}

The second lemma, which by \S.\,181 still remains to be proved, is the
following:

2. \gesperrt{If $\varepsilon_n$ is a system of conjugate numerical units
of the field $\Omega_m$ which satisfies the condition that}

\weq{1}{\varepsilon_n \varepsilon_1^{-n} = \mathfrak{E}^m}

\gesperrt{is the $m$th power of a unit in $\Omega_m$, then the
$\varepsilon_n$ are themselves $m$th powers of units in $\Omega_m$,
multiplied by a root of unity of the form $r^{nk}$.}

This proof too is entirely elementary in the case of an odd $m$. We
therefore consider this case here first.

By \S.\,178, 1. we can represent every unit of the field $\Omega_m$ in
the following form:

\weq{2}{\varepsilon = \pm r^k \mathfrak{E}(r), \quad \varepsilon_n = \pm r^{nk} \mathfrak{E}(r^n),}

where $\mathfrak{E}(r)$ is a real unit of the field $\Omega_m$.
Accordingly $\mathfrak{E}(r)$ satisfies the condition

\weq{3}{\mathfrak{E}(r) = \mathfrak{E}(r^{-1}),}

from which, by (2), there results:

\weq{4}{\varepsilon_1 \varepsilon_{-1} = \mathfrak{E}(r)^2;}

but from (1) one finds, on putting $n = -1$,

\[
\varepsilon_1 \varepsilon_{-1} = \mathfrak{E}^m,
\]

that is, the left side of (4) is the $m$th power of a unit in
$\Omega_m$, and therefore

\weq{5}{\mathfrak{E}(r)^2 = \mathfrak{E}^m}

is likewise the $m$th power of such a unit.

Since $m$ is now odd, the integral rational number $x$ can be so
determined that

\weq{6}{2 x + m = 1}

holds, and from this it follows that

\[
\mathfrak{E}(r) = \mathfrak{E}(r)^{2x} \mathfrak{E}(r)^m.
\]
%CONT>>
